% Осень26!A34; нумерация внутри столбца без пустых строк.
% Источник: лекция 7, §3 (строка 65)
\ticket{31}{Критерий плотности}
\label{ticket:31}

\begin{theorem}
Пусть $\mu,\nu$ --- меры на $\sigma$-алгебре $\mathcal A$ подмножеств $X$, а $\omega:X\to[0,+\infty]$ измерима.
Функция $\omega$ является плотностью $\nu$ относительно $\mu$ тогда и только тогда, когда для любого непустого $A\in\mathcal A$
\[
  \mu(A)\inf_{x\in A}\omega(x)
  \le\nu(A)\le
  \mu(A)\sup_{x\in A}\omega(x).
\]
Здесь $0\cdot\infty=\infty\cdot0=0$.
\end{theorem}

\begin{proof}
Необходимость следует из оценки интеграла снизу и сверху постоянными на $A$.

Докажем достаточность.
На $A_0=A\cap\{\omega=0\}$ верхняя оценка даёт $\nu(A_0)=0=\int_{A_0}\omega\,d\mu$.
На $A_\infty=A\cap\{\omega=\infty\}$ при $\mu(A_\infty)=0$ обе величины равны нулю, а при $\mu(A_\infty)>0$ обе бесконечны по нижней оценке.

Остаётся $A'=A\cap\{0<\omega<\infty\}$.
Для $q\in(0,1)$ положим
\[
  A_j=\{x\in A':q^j\le\omega(x)<q^{j-1}\},\qquad j\in\mathbb Z.
\]
Это счётное разбиение $A'$.
Из условия и монотонности интеграла следуют оценки
\[
  q^j\mu(A_j)\le\nu(A_j)\le q^{j-1}\mu(A_j),
\]
\[
  q^j\mu(A_j)\le I_j:=\int_{A_j}\omega\,d\mu
  \le q^{j-1}\mu(A_j).
\]
Поэтому $qI_j\le\nu(A_j)\le q^{-1}I_j$.
Суммируя по $j$ и используя счётную аддитивность меры и интеграла по множеству, получаем
\[
  q\int_{A'}\omega\,d\mu\le\nu(A')
  \le q^{-1}\int_{A'}\omega\,d\mu.
\]
При конечном интеграле устремляем $q\to1$; при бесконечном нижняя оценка сразу даёт $\nu(A')=\infty$.
Следовательно, $\nu(A')=\int_{A'}\omega\,d\mu$.
Добавляя части $A_0,A_\infty$, получаем равенство на всём $A$.
\end{proof}
